\chapter{Testování}
\label{ch:testovani}

Tato kapitola popisuje průběh a výsledky uživatelského testování aplikace Foody.
Cílem testování bylo ověřit použitelnost implementované aplikace na reálných uživatelích a identifikovat případné problémy v~uživatelském rozhraní a interakčních tocích.
V~první části kapitoly je popsána zvolená metodika, výběr participantů, sada testovacích úloh a použité dotazníky.
Následně jsou prezentovány výsledky plnění jednotlivých úloh, výstupy standardizovaných dotazníků a identifikované problémy s~doporučeními pro další iteraci vývoje.

\section{Metodika testování}
\label{sec:metodikaTestovani}

V~rámci testování aplikace Foody byla zvolena kombinace metody \emph{think aloud} a úlohového testování (\emph{task-based usability testing}).
Tato kombinace umožňuje získat jak kvantitativní data o~úspěšnosti a rychlosti plnění úloh, tak kvalitativní poznatky o~kognitivních procesech participantů při interakci s~aplikací.
Následující podkapitoly popisují výběr participantů, prostředí testování, sadu testovacích úloh a použité dotazníky a metriky.

\subsection{Výběr a charakteristika participantů}
\label{subsec:vyberParticipantu}

Pro uživatelské testování byli vybráni celkem čtyři participanti, kteří představují různé uživatelské profily z~hlediska zkušeností se sledováním stravy a technologické zdatnosti.
Tento počet je v~souladu s~doporučeními Nielsena, podle kterých pět participantů odhalí přibližně 80~\% problémů použitelnosti, přičemž již čtyři participanti pokrývají většinu kritických nedostatků~\cite{nielsen2000fiveUsers}.
Participanti byli vybráni s~ohledem na variabilitu věku, pohlaví a předchozích zkušeností s~aplikacemi pro sledování kalorického příjmu, čímž bylo zajištěno pokrytí rozdílných uživatelských perspektiv.

% TODO: doplnit profily participantů P1--P4 (jedná se o stejné participanty jako v uživatelské studii v sekci 3.2.2).
% Pro každého uvést primární cíl, zkušenosti s calorie tracking aplikacemi a technologickou zdatnost (2--4 věty).
% Vzor: „P1 (žena, 28 let) se aktivně věnuje fitness a sledování stravy. Dříve používala MyFitnessPal po dobu šesti měsíců,
% kterou opustila kvůli časové náročnosti manuálního zápisu. Od testované aplikace očekává rychlejší záznam díky AI rozpoznávání."

\subsection{Prostředí a průběh testovací relace}
\label{subsec:prostrediRelace}

Testování probíhalo formou individuálních relací v~kontrolovaném prostředí, kdy každý participant pracoval s~aplikací na reálném mobilním zařízení.
V~průběhu testování byla použita metoda souběžného protokolu \emph{think aloud}, při které participanti verbalizovali své myšlenky, dojmy a pochybnosti během plnění úloh~\cite{nielsen1994thinkAloud}.
Na začátku každé relace byl participant seznámen s~účelem testování, podepsal informovaný souhlas a obdržel instrukce k~verbalizaci myšlenek.

Každá testovací relace měla odhadovanou délku 75 až 90~minut a sestávala ze čtyř fází, jejichž přehled je uveden v~tabulce~\ref{tab:strukturaRelace}.
V~úvodní fázi (přibližně 5~minut) byl participant seznámen s~průběhem testu a instrukcemi k~metodě \emph{think aloud}.
Hlavní fáze testování (60 až 70~minut) zahrnovala plnění třinácti testovacích úloh, přičemž po dokončení každé úlohy participant vyplnil dotazník SEQ.
Třetí fáze (přibližně 10~minut) byla věnována vyplnění standardizovaných dotazníků SUS a UEQ-S hodnotících celkovou zkušenost s~aplikací.
Závěrečná fáze (5 až 10~minut) sestávala z~polostrukturovaného debriefingu s~šesti otevřenými otázkami zaměřenými na celkový dojem, pozitivní a negativní aspekty aplikace, srovnání s~konkurenčními řešeními a ochotu k~pravidelnému používání.

\begin{table}[ht]
	\centering
	\caption{Struktura testovací relace.}
	\label{tab:strukturaRelace}
	\begin{tabular}{lll}
		\toprule
		\textbf{Fáze} & \textbf{Délka} & \textbf{Obsah} \\
		\midrule
		Úvod & 5~min & Představení, souhlas, instrukce k~\emph{think aloud} \\
		Testovací úlohy & 60--70~min & Plnění úloh T1--T13 s~SEQ po každé úloze \\
		Dotazníky & 10~min & SUS a UEQ-S \\
		Debriefing & 5--10~min & Otevřené otázky, celkový dojem \\
		\bottomrule
	\end{tabular}
\end{table}

Na základě přípravy testovacího prostředí byla na zařízení nainstalována čistá verze aplikace bez předchozích dat, čímž byl simulován stav prvního spuštění.
Do galerie zařízení byly nahrány připravené fotografie jídel pro úlohu T3, na stůl bylo připraveno reálné jídlo pro úlohu T2 a balený produkt s~čárovým kódem pro úlohu T6.
Průběh testování byl zaznamenáván formou záznamu obrazovky a zvukového záznamu pro následnou analýzu verbalizací.

\subsection{Testovací úlohy}
\label{subsec:testovaciUlohy}

Testovací úlohy byly navrženy tak, aby pokrývaly klíčové funkční požadavky a případy užití definované v~kapitole~\ref{ch:navrh} (viz sekce~\ref{sec:uzivatelskaStudie}).
Úlohy jsou seřazeny v~pořadí, které simuluje přirozený průchod aplikací od prvního spuštění přes různé způsoby záznamu jídla až po přehledy a export dat.
Celkem bylo definováno třináct testovacích úloh, jejichž přehled včetně vazby na funkční požadavky a případy užití je uveden v~tabulce~\ref{tab:ulohy}.

\begin{table}[ht]
	\centering
	\caption{Přehled testovacích úloh s~vazbou na funkční požadavky a případy užití.}
	\label{tab:ulohy}
	\begin{tabular}{cll}
		\toprule
		\textbf{Úloha} & \textbf{Název} & \textbf{Pokryté FR / UC} \\
		\midrule
		T1 & Onboarding a nastavení profilu & FR-03, FR-04, FR-20, UC07 \\
		T2 & Záznam jídla z~fotografie & FR-06, FR-07, FR-08, FR-09, UC01 \\
		T3 & Import fotografie z~galerie & FR-12 \\
		T4 & Oprava výsledku rozpoznání a re-analýza & FR-11, FR-13, UC02 \\
		T5 & Záznam jídla hlasem & FR-14 \\
		T6 & Skenování čárového kódu & FR-16, UC04 \\
		T7 & Ruční přidání záznamu & FR-01, FR-15, UC03 \\
		T8 & Oblíbené a duplikace jídla & FR-17, FR-18, UC05 \\
		T9 & Záznam cvičení hlasem & --- \\
		T10 & Zaznamenání váhy & --- \\
		T11 & Denní přehled a smazání záznamu & FR-02, FR-05, FR-22 \\
		T12 & Týdenní přehled a Ask AI & FR-24, FR-27, UC08 \\
		T13 & Export dat & FR-25 \\
		\bottomrule
	\end{tabular}
\end{table}

Úloha T1 zahrnovala kompletní onboarding flow aplikace, v~rámci kterého participant zadal základní údaje (váha, výška, věk), nastavil kalorický cíl a zvolil případné dietní preference.
Úlohy T2 až T6 pokrývaly jednotlivé vstupní modality pro záznam jídla, konkrétně fotografii z~kamery (T2), import z~galerie (T3), opravu a re-analýzu výsledku rozpoznání (T4), hlasový vstup (T5) a skenování čárového kódu (T6).
Úloha T7 testovala ruční přidání záznamu bez využití umělé inteligence, zatímco úloha T8 ověřovala práci s~oblíbenými položkami a funkcí kopírování jídla na jiný den.

V~rámci úloh T9 a T10 bylo testováno rozšíření aplikace nad rámec původní specifikace, konkrétně záznam cvičení hlasovým vstupem a zaznamenání tělesné váhy.
Úlohy T11 až T13 se zaměřovaly na přehledové a analytické funkce aplikace, včetně denního přehledu se smazáním záznamu (T11), týdenního přehledu s~funkcí Ask AI pro dotazy v~přirozeném jazyce (T12) a exportu dat (T13).

Z~celkového počtu třiceti funkčních požadavků definovaných v~kapitole~\ref{ch:navrh} je testovacími úlohami přímo pokryto dvacet dva požadavků, což představuje 73~\% specifikace.
Všech osm případů užití (UC01 až UC08) je pokryto alespoň jednou úlohou.
Zbývajících osm funkčních požadavků nebylo zahrnuto do testování z~důvodů, které jsou shrnuty v~tabulce~\ref{tab:nepokrytaFR}.

\begin{table}[ht]
	\centering
	\caption{Funkční požadavky nepokryté testovacími úlohami s~uvedením důvodu.}
	\label{tab:nepokrytaFR}
	\small
	\begin{tabular}{p{1.1cm}p{4.5cm}p{8cm}}
		\toprule
		\textbf{FR} & \textbf{Název} & \textbf{Důvod vynechání} \\
		\midrule
		FR-10 & Rozlišení chyb rozpoznávání a aplikace & Nelze spolehlivě vyvolat v~kontrolovaném prostředí \\
		FR-19 & Našeptávání názvů & Implicitně pokryto vyhledáváním v~úlohách T7 a T8 \\
		FR-21 & Porušení diet v~kalendáři & Vyžaduje historická data s~porušeními, obtížné simulovat \\
		FR-23 & Nastavení integrace výdeje & Pouze přepínač v~nastavení, nízká interakční hodnota \\
		FR-26 & Offline tolerance & Vyžaduje odpojení od sítě během testu \\
		FR-28 & Motivační souhrn & Vyžaduje dlouhodobý cyklus (denní, týdenní, měsíční) \\
		FR-29 & Konfigurovatelné notifikace & Efekt vyžaduje časový odstup \\
		FR-30 & Zobrazení a skrytí pokročilých funkcí & Pouze přepínače v~nastavení, nízká interakční hodnota \\
		\bottomrule
	\end{tabular}
\end{table}

Během plnění vybraných úloh bylo současně pozorováním ověřováno plnění nefunkčních požadavků.
V~rámci úlohy T2 byla měřena latence rozpoznávání umělou inteligencí (NFR-02, limit 20~s) časem od odeslání fotografie do zobrazení výsledku.
Počet kroků potřebných pro zápis jídla (NFR-03, limit 6~kroků) byl zaznamenáván při úlohách T2, T5, T6 a T7.
Čas zápisu (NFR-04) byl měřen pro nový záznam v~úloze T2 (limit 5~minut) a pro opakovaný záznam v~úloze T8 (limit 1~minuta).

\subsection{Použité dotazníky a metriky}
\label{subsec:dotazniky}

Pro kvantitativní hodnocení použitelnosti a uživatelské zkušenosti byly v~rámci testování použity tři standardizované dotazníky.
Výběr nástrojů pokrývá hodnocení na úrovni jednotlivých úloh (SEQ), celkovou použitelnost systému (SUS) a uživatelskou zkušenost v~pragmatické i hedonické dimenzi (UEQ-S).
Společné použití těchto tří dotazníků umožňuje triangulaci výsledků a poskytuje komplexní pohled na kvalitu interakce uživatele s~aplikací.

Bezprostředně po dokončení každé testovací úlohy participant vyplnil dotazník \emph{Single Ease Question} (SEQ), který sestává z~jediné otázky hodnotící subjektivní obtížnost úlohy na sedmibodové škále, kde 1 znamená „velmi obtížná" a 7 „velmi snadná"~\cite{sauro2009seq}.
Výhodou SEQ je minimální zatížení participanta a možnost okamžitého zachycení dojmu z~právě dokončené úlohy, aniž by byl narušen průběh testování.
Na základě dostupných benchmarků je za dobré hodnocení považován průměr SEQ skóre 5,5 a vyšší.

Po dokončení všech testovacích úloh participant vyplnil dotazník \emph{System Usability Scale} (SUS), který představuje standardizovaný nástroj pro hodnocení celkové použitelnosti systému~\cite{brooke1996sus}.
Dotazník SUS sestává z~deseti položek hodnocených na pětibodové Likertově škále (1~--~rozhodně nesouhlasím, 5~--~rozhodně souhlasím), přičemž liché položky jsou formulovány pozitivně a sudé negativně.
Výsledné skóre se pohybuje v~rozsahu 0 až 100~bodů a vypočítá se podle vzorce~(\ref{eq:sus}).
V~kontextu dostupných benchmarků odpovídá průměrné SUS skóre napříč studiemi přibližně 68~bodům, přičemž skóre nad 80,3~bodu řadí systém mezi horních 10~\% hodnocených aplikací~\cite{bangor2008susAdjective}.

\begin{equation}
	\mathrm{SUS} = \bigl((S_L - 5) + (25 - S_S)\bigr) \cdot 2{,}5,
	\label{eq:sus}
\end{equation}

\noindent kde $S_L$ je součet odpovědí na lichých položkách (1, 3, 5, 7, 9) a $S_S$ je součet odpovědí na sudých položkách (2, 4, 6, 8, 10).

Jako doplňkový nástroj pro hodnocení uživatelské zkušenosti byl použit dotazník \emph{User Experience Questionnaire Short} (UEQ-S), který měří pragmatickou a hedonickou kvalitu systému prostřednictvím osmi položek~\cite{schrepp2017ueqs}.
Každá položka má formu sémantického diferenciálu na sedmibodové škále ($-3$ až $+3$), přičemž záporný pól je vždy umístěn na levé straně.
Položky 1 až 4 (bránící/podporující, složitý/jednoduchý, neefektivní/efektivní, matoucí/jasný) měří pragmatickou kvalitu a položky 5 až 8 (nudný/vzrušující, nezajímavý/zajímavý, obvyklý/vynalézavý, tradiční/moderní) měří hedonickou kvalitu.
Pragmatická kvalita se vypočítá jako průměr položek 1 až 4, hedonická kvalita jako průměr položek 5 až 8 a celkové skóre jako průměr všech osmi položek.
Byla použita oficiální česká verze dotazníku validovaná a dostupná na stránkách projektu UEQ.

\section{Výsledky testování}
\label{sec:vysledkyTestovani}

Následující podkapitoly prezentují výsledky uživatelského testování v~členění na kvantitativní údaje o~plnění jednotlivých úloh, hodnocení SEQ a výstupy standardizovaných dotazníků SUS a UEQ-S.
Výsledky jsou doplněny o~kvalitativní poznatky získané analýzou verbalizací participantů v~průběhu metody \emph{think aloud} a z~odpovědí debriefingových rozhovorů.

\subsection{Plnění úloh a hodnocení SEQ}
\label{subsec:plneniUloh}

% TODO: pro každou úlohu T1--T13 uvést v průvodním textu:
%   --- míra dokončení (kolik ze 4 participantů úspěšně dokončilo)
%   --- průměrný čas plnění
%   --- počet kritických chyb (participant nedokázal pokračovat bez pomoci moderátora)
%   --- průměrné SEQ skóre
%   --- reprezentativní pozorování z think aloud (parafráze, ne přímé citáty)
% Doporučený rozsah: 1--2 odstavce na úlohu; výrazné nálezy rozepsat (T2, T4, T5, T8),
% bezproblémové úlohy lze sloučit do souhrnného odstavce.

\begin{table}[ht]
	\centering
	\caption{Výsledky plnění testovacích úloh a hodnocení SEQ.}
	\label{tab:vysledkyUloh}
	\begin{tabular}{cccccc}
		\toprule
		\textbf{Úloha} & \textbf{Dokončení (n/4)} & \textbf{$\varnothing$ čas (s)} & \textbf{Kritické chyby} & \textbf{$\varnothing$ SEQ (1--7)} \\
		\midrule
		T1 & TODO & TODO & TODO & TODO \\
		T2 & TODO & TODO & TODO & TODO \\
		T3 & TODO & TODO & TODO & TODO \\
		T4 & TODO & TODO & TODO & TODO \\
		T5 & TODO & TODO & TODO & TODO \\
		T6 & TODO & TODO & TODO & TODO \\
		T7 & TODO & TODO & TODO & TODO \\
		T8 & TODO & TODO & TODO & TODO \\
		T9 & TODO & TODO & TODO & TODO \\
		T10 & TODO & TODO & TODO & TODO \\
		T11 & TODO & TODO & TODO & TODO \\
		T12 & TODO & TODO & TODO & TODO \\
		T13 & TODO & TODO & TODO & TODO \\
		\bottomrule
	\end{tabular}
\end{table}

% TODO: souhrnná interpretace výsledků z tabulky:
%   --- které úlohy měly nejvyšší a nejnižší SEQ skóre
%   --- kde se vyskytly kritické chyby a co je způsobilo
%   --- celkový průměr SEQ napříč úlohami a srovnání s benchmarkem 5,5
%   --- výsledky ověření NFR-02 (latence), NFR-03 (počet kroků), NFR-04 (čas zápisu)
%   --- shrnutí, které toky fungují dobře a které vyžadují pozornost.

\subsection{Výsledky System Usability Scale}
\label{subsec:vysledkySus}

% TODO: interpretace SUS výsledků. Vzor:
% „Průměrné SUS skóre aplikace Foody činilo X,X bodu (SD = X,X), což ji řadí do kategorie X
% podle adjektivní stupnice Bangor a kol. Nejvyšší hodnocení získala položka č. X (průměr X,X),
% která se týká Y. Naopak nejnižší hodnocení obdržela položka č. X (průměr X,X), což naznačuje,
% že participanti vnímali Z jako problematický aspekt. Celkově lze konstatovat, že dosažené
% skóre X,X odpovídá hodnocení 'X' a řadí aplikaci nad/pod průměr SUS benchmarku 68 bodů."

\begin{table}[ht]
	\centering
	\caption{Výsledky dotazníku SUS pro jednotlivé participanty.}
	\label{tab:vysledkySus}
	\begin{tabular}{lc}
		\toprule
		\textbf{Participant} & \textbf{SUS skóre} \\
		\midrule
		P1 & TODO \\
		P2 & TODO \\
		P3 & TODO \\
		P4 & TODO \\
		\midrule
		\textbf{Průměr} & TODO \\
		\textbf{Směrodatná odchylka} & TODO \\
		\bottomrule
	\end{tabular}
\end{table}

\subsection{Výsledky User Experience Questionnaire Short}
\label{subsec:vysledkyUeqs}

% TODO: interpretace UEQ-S výsledků. Vzor:
% „Pragmatická kvalita aplikace dosáhla průměrného skóre X,X, což podle benchmarků UEQ
% odpovídá hodnocení 'X'. Hedonická kvalita dosáhla skóre X,X. Z jednotlivých položek získalo
% nejvyšší hodnocení sémantické páry X/X (průměr X,X), zatímco nejnižší hodnocení obdržely
% páry X/X (průměr X,X). Výsledky naznačují, že participanti vnímají aplikaci jako pragmaticky X,
% avšak v hedonické dimenzi existuje prostor pro zlepšení v oblasti Y."

\begin{table}[ht]
	\centering
	\caption{Výsledky dotazníku UEQ-S.}
	\label{tab:vysledkyUeqs}
	\begin{tabular}{clc}
		\toprule
		\textbf{\#} & \textbf{Položka} & \textbf{Průměr} \\
		\midrule
		1 & bránící / podporující & TODO \\
		2 & složitý / jednoduchý & TODO \\
		3 & neefektivní / efektivní & TODO \\
		4 & matoucí / jasný & TODO \\
		5 & nudný / vzrušující & TODO \\
		6 & nezajímavý / zajímavý & TODO \\
		7 & obvyklý / vynalézavý & TODO \\
		8 & tradiční / moderní & TODO \\
		\midrule
		--- & \textbf{Pragmatická kvalita} & TODO \\
		--- & \textbf{Hedonická kvalita} & TODO \\
		--- & \textbf{Celkové skóre} & TODO \\
		\bottomrule
	\end{tabular}
\end{table}

\section{Identifikované problémy}
\label{sec:identifikovaneProblemy}

Na základě analýzy pozorování z~metody \emph{think aloud}, výsledků dotazníků a debriefingových rozhovorů byly identifikovány problémy použitelnosti.
Každý problém je klasifikován podle závažnosti na jedné ze tří úrovní: kritický (participant nebyl schopen úlohu dokončit bez pomoci moderátora), závažný (participant úlohu dokončil, ale s~výrazným zdržením nebo zmateností) a kosmetický (participant si všiml nedostatku, avšak ten neovlivnil plnění úlohy).
Přehled identifikovaných problémů je uveden v~tabulce~\ref{tab:problemy}.

\begin{table}[ht]
	\centering
	\caption{Identifikované problémy použitelnosti s~klasifikací závažnosti.}
	\label{tab:problemy}
	\begin{tabular}{p{1cm}p{6cm}p{2.4cm}p{1.5cm}p{1.5cm}}
		\toprule
		\textbf{\#} & \textbf{Problém} & \textbf{Závažnost} & \textbf{Úloha} & \textbf{FR / NFR} \\
		\midrule
		% TODO: doplnit identifikované problémy (5--12 dle skutečných nálezů).
		% Vzor:
		% P1 & Nízká viditelnost indikátoru spolehlivosti rozpoznání & Závažný & T2 & FR-08 \\
		% P2 & Obtížná nalezitelnost funkce kopírování jídla & Závažný & T8 & FR-18 \\
		TODO & TODO & TODO & TODO & TODO \\
		\bottomrule
	\end{tabular}
\end{table}

% TODO: za tabulkou popsat každý problém 2--3 větami. Vzor:
% „V rámci úlohy T2 dva ze čtyř participantů nepovšimli barevného indikátoru spolehlivosti
% výsledku (confidence badge). Participanti tak neměli povědomí o míře spolehlivosti rozpoznání,
% což snižuje efektivitu funkčního požadavku FR-08. Problém se pravděpodobně týká vizuálního
% návrhu badge, který splývá s okolním rozhraním."

\section{Doporučení pro další iteraci}
\label{sec:doporuceni}

Na základě identifikovaných problémů a výsledků standardizovaných dotazníků lze formulovat doporučení pro další vývojovou iteraci aplikace.
Doporučení jsou řazena podle priority, která vychází ze závažnosti příslušného problému a počtu dotčených participantů.
Každé doporučení se váže ke konkrétnímu identifikovanému problému z~kapitoly~\ref{sec:identifikovaneProblemy}.

% TODO: doporučení organizovaná do odstavců (4--8 dle počtu závažných a kritických problémů).
% Každé doporučení obsahuje:
%   --- odkaz na číslo problému z 5.3
%   --- konkrétní návrh řešení
%   --- dotčené FR/NFR
%   --- očekávaný dopad na použitelnost
% Vzor:
% „V návaznosti na problém P1 (nízká viditelnost indikátoru spolehlivosti) je doporučeno zvýšit
% vizuální prominenci confidence badge zvětšením jeho rozměrů a přidáním textového popisku
% vedle barevného indikátoru. Tato úprava by měla zvýšit povědomí uživatelů o míře spolehlivosti
% výsledku v souladu s požadavkem FR-08."

% TODO: závěrečný odstavec kapitoly. Vzor:
% „Z výsledků testování vyplývá, že aplikace Foody dosáhla X,X bodů v dotazníku SUS a průměrného
% SEQ skóre X,X napříč všemi úlohami. Pragmatická kvalita dle UEQ-S byla hodnocena jako 'X',
% což naznačuje, že základní interakční toky aplikace jsou pro uživatele srozumitelné a efektivní.
% Identifikované problémy se soustředí především na X, přičemž navržená doporučení jsou
% realizovatelná v rámci další vývojové iterace bez zásadních architektonických změn. Celkově lze
% konstatovat, že výsledky testování potvrzují životaschopnost zvoleného přístupu k rozpoznávání
% potravin pomocí umělé inteligence s transparentní indikací nejistoty, avšak poukazují na konkrétní
% oblasti, kde je třeba zlepšit viditelnost klíčových prvků rozhraní a nalezitelnost sekundárních funkcí."
